\section{Exact pair-key geometry}

Throughout the following finite-field arithmetic sections, p is an odd prime,
$H\leq\mathbb F_p^\times$ has odd prime order K, and
$e_p(t)=e^{2\pi i t/p}$. Fixed-degree row presentations impose $K\geq7$
where stated. These hypotheses are distinct from the composite Pell moduli
in the later quantum sections.

For two admissible position pairs $(x_i,y_i)\in(\mathbb F_p\setminus\{0,1\})^2$,
write their keys as $(u_i,v_i,d_i)$. Let $B(t)$ count the ordered quadruples
satisfying $u_1u_2\in H$, $v_1v_2\in H$, and $d_1+d_2=t$; let $O(t)$ count
those for which both position pairs are off diagonal. Write $V_4$ for the
set of all such quadruples without fixing t. These definitions include
multiplicities in the collapsed diagonal fibre.

For $x,y\in\Fp\setminus\{0,1\}$ define
\begin{equation}\label{eq:keydef}
 u=\frac{x}{y},\qquad
 v=\frac{x-1}{y-1},\qquad
 d=x-y.
\end{equation}

\begin{proposition}[Pair-key bijection]\label{prop:key}
If $x\ne y$, then $u,v\in\Fp^\times\setminus\{1\}$, $u\ne v$, and
\begin{equation}\label{eq:inverse}
 y=\frac{1-v}{u-v},\qquad
 x=\frac{u(1-v)}{u-v},\qquad
 d=\frac{(u-1)(v-1)}{v-u}.
\end{equation}
Conversely every such $(u,v)$ gives a unique admissible off-diagonal pair.
The diagonal collapses to $(u,v,d)=(1,1,0)$ with multiplicity $p-2$.
\end{proposition}

\begin{proof}
The equations $x=uy$ and $x-1=v(y-1)$ give $(u-v)y=1-v$, hence
\eqref{eq:inverse}.  The diagonal statement is immediate.  Conversely, given
$u,v\in\Fp^\times\setminus\{1\}$ with $u\ne v$, define $x,y$ by
\eqref{eq:inverse}.  Then $y\ne0$ because $v\ne1$, and $y\ne1$ because $u\ne1$;
$x\ne0$ because $u\ne0$ and $v\ne1$; and $x\ne1$, since $x=1$ would give
$u(1-v)=u-v$, hence $uv=v$ and $u=1$.  Also $x=uy$ with $u\ne1$ and $y\ne0$
gives $x\ne y$, and substituting into \eqref{eq:keydef} returns $(u,v)$.
\end{proof}

Under the subgroup constraints $u,v\in H$, there are exactly
$(K-1)(K-2)$ admissible off-diagonal ordered pairs.  If $O(t)$ denotes the
off/off contribution and
\[
C_H(t)=\#\left\{(u,v)\in(H\setminus\{1\})^2:
 u\ne v,\ \frac{(u-1)(v-1)}{v-u}=t\right\},
\]
then exact bookkeeping gives
\begin{equation}\label{eq:Bdecomp}
B(t)=O(t)+2(p-2)C_H(t)+(p-2)^2\mathbf1_{t=0}.
\end{equation}
The factor two is essential: diagonal/off and off/diagonal are distinct ordered
classes.  The total count is
\begin{equation}\label{eq:V4}
|V_4|=(K^2+1)p^2-(6K^2+3K)p+13K^2-3K.
\end{equation}
If $\Sigma_4(q)=\sum_tB(t)e_p(qt)$, $V=\sum_tB(t)$ and $Z=B(0)$, finite
Fourier orthogonality gives
\begin{equation}\label{eq:meanvar}
\mathbb E_{q\ne0}\Sigma_4(q)=\frac{pZ-V}{p-1},
\qquad
\Var_{q\ne0}\Sigma_4(q)=p\Var_{t\ne0}B(t).
\end{equation}
Here, and throughout, $\Var$ of a complex-valued quantity denotes
$\mathbb E|\,\cdot\,-\mathbb E(\cdot)|^2$ over the indicated finite index
set, with $\mathbb E$ the uniform average; for real-valued quantities this is
the ordinary variance.
No probabilistic independence assumption enters this reduction.

\section{Cayley energy and a uniform bound}

Put
\[
c(h)=\frac{h+1}{h-1},\qquad A=c(H\setminus\{1\}).
\]
For $u\ne v$,
\begin{equation}\label{eq:cayley}
\frac{(u-1)(v-1)}{v-u}=\frac{2}{c(u)-c(v)},
\qquad c(h^{-1})=-c(h).
\end{equation}
Thus $A=-A$ and the nonzero multiplicities of $C_H$ are additive-difference
multiplicities of $A$ up to the permutation $x\mapsto2/x$.  Hence
\begin{equation}\label{eq:energyCH}
\sum_tC_H(t)^2=E_+(A)-|A|^2,
\qquad
E_+(A)=\sum_xr_{A-A}(x)^2.
\end{equation}
Identity and central inversion force the nonzero baseline
\begin{equation}\label{eq:forced}
(K-1)(2K-5).
\end{equation}
For the unnormalized Fourier transform,
\begin{equation}\label{eq:centered}
E_+(A)=\frac{(K-1)^4}{p}
 +\frac1p\sum_{q\ne0}|\widehat{1_A}(q)|^4.
\end{equation}
The centered identity \eqref{eq:centered} separates the unavoidable
zero-mode term $(K-1)^4/p$.  In the regime where that term is $O(K^2)$ -- in
particular in the intended range $K=o(p^{1/4})$ -- the forced baseline
\eqref{eq:forced} sets the quadratic scale, and the sharp pointwise target
(open, not a theorem) is
\begin{equation}\label{eq:sharp-target}
E_+(A)=\frac{(K-1)^4}{p}+O(K^{2+o(1)}).
\end{equation}

For $r_A(t)=\#\{(a,b)\in A^2:a+b=t\}$, clearing Cayley denominators gives
\[
(2-t)xy+t(x+y)-(2+t)=0,
\]
an absolutely irreducible curve in odd characteristic: when $t\ne2$ the
bilinear coefficient determinant is $-4\ne0$, and when $t=2$ it is a
nonconstant linear polynomial. Its bidegree is $(1,1)$ in either case.
For $t\ne0,-2$ the constant term is nonzero and $\deg P(x,0)=1$, exactly
the hypotheses needed for the subgroup-curve input.

\begin{theorem}[Direct pointwise Cayley bound]\label{thm:direct}
If
\[
100<K<\frac13p^{3/4},
\]
then
\begin{equation}\label{eq:direct}
E_+(A)<34K^{8/3}.
\end{equation}
\end{theorem}

\begin{proof}
The difference energy $E_+(A)$ of \eqref{eq:energyCH} equals the
sum-representation energy: $\sum_xr_{A-A}(x)^2$ and $\sum_tr_A(t)^2$ both
count $\{(a,b,c,d)\in A^4:a+b=c+d\}$, after relabelling the second and
fourth variables.
Theorem~2 of Vyugin--Makarychev \cite{VM2015,MV2019}, in the source form for
bidegree $(m,n)=(1,1)$, applies to an absolutely irreducible $P$ with
$P(0,0)\ne0$ and $\deg P(x,0)\ge1$ under the size condition
$100(mn)^{3/2}<K<\frac13p^{3/4}$ -- for $m=n=1$ exactly the hypothesis above --
and gives $r_A(t)\le16mn^2(m+n)K^{2/3}=32K^{2/3}$.  On the displayed
polynomial, $P_t(0,0)=-(2+t)$ and $P_t(x,0)=tx-(2+t)$, so the two degree
hypotheses fail only at $t=-2$ and at $t=0$ respectively.  The fibre $t=2$
is not exceptional: $P_2(x,y)=2(x+y-2)$ is still absolutely irreducible of
bidegree $(1,1)$, with $P_2(0,0)=-4$ and $\deg P_2(x,0)=1$, so the same
bound $r_A(2)\le32K^{2/3}$ holds there.  The exceptional fibres satisfy
$r_A(0)=K-1$ and $r_A(-2)\le K-1$.  Since
$\sum_tr_A(t)=(K-1)^2$,
\[
E_+(A)
\le2(K-1)^2+32K^{2/3}\sum_{t\ne0,-2}r_A(t)
<34K^{8/3}.
\]
\end{proof}

\section{Prime-order rigidity, square resultants and discriminants}

Let $\zeta=e^{2\pi i/K}$.  Then
\[
c(\zeta^a)=-i\cot\frac{\pi a}{K}.
\]
Chowla's prime-modulus cotangent independence \cite{Chowla1970,Wang1982}
identifies the characteristic-zero Cayley set with a rational cross-polytope.

\begin{theorem}[Prime-order additive rigidity]\label{thm:add-rigidity}
Every characteristic-zero equal-difference collision among
$\{c(\zeta^a):a\in\mathbf F_K^*\}$ is identity or central inversion-swap.
\end{theorem}

\begin{proof}
Write $b_a=\cot(\pi a/K)$ for $1\le a\le m=(K-1)/2$.  By Chowla's theorem
\cite{Chowla1970,Wang1982} the numbers $b_a$ are linearly independent over
$\mathbf Q$.  Since $c(\zeta^a)=-i\cot(\pi a/K)$ and
$\cot(\pi(K-a)/K)=-\cot(\pi a/K)$, every value $c(\zeta^a)$ is $-i$ times a
signed vertex $\pm b_a$; under $b_a\leftrightarrow e_a$ the Cayley values are
the vertices of the rational cross-polytope $\{\pm e_a\}$.  An
equal-difference collision $c(\zeta^i)-c(\zeta^j)=c(\zeta^r)-c(\zeta^s)$,
$i\ne j$, $r\ne s$, therefore gives an identity $x-y=u-v$ among signed
vertices with $x\ne y$ and $u\ne v$, i.e. $x+v=u+y$.  If $x+v=0$, then
$v=-x$ and $u=-y$: the central inversion-swap $(r,s)=(-j,-i)$.  Otherwise
$x+v\ne0$ is a sum of two signed basis vectors, and such a sum determines
its unordered pair of summands uniquely -- including the diagonal case
$v=x$, where the sum is $\pm2e_a$ and both summands are forced.  Hence
$\{x,v\}=\{u,y\}$; the alternative $u=y$, $v=x$ would give $x-y=y-x$,
forcing $x=y$, a contradiction.  Thus $u=x$ and $v=y$: the identity
collision.
\end{proof}

For every admissible cleared collision polynomial $P(X)$, direct algebra gives
$P(X)=(X-1)^3Q(X)$ and a reciprocal relation for $Q$.  Passing to the maximal
real subfield yields:

\begin{theorem}[Square-resultant theorem]\label{thm:square-res}
\begin{equation}\label{eq:square-res}
\Res(\Phi_K,P)=K^3M^2,\qquad M\in\Z.
\end{equation}
For a non-forced characteristic-zero collision, $M\ne0$.
\end{theorem}

\begin{proof}
Here the cleared equal-difference polynomial can be defined explicitly.
For $1\leq a_i<K$ and $(\epsilon_1,\epsilon_2,\epsilon_3,\epsilon_4)
=(1,-1,-1,1)$, put
\[
 P(X)=\sum_{i=1}^4\epsilon_i(X^{a_i}+1)
                         \prod_{j\ne i}(X^{a_j}-1),\qquad s=\sum_i a_i.
\]
Every term contains three factors vanishing at one, so
$P=(X-1)^3Q$ with $Q\in\mathbb Z[X]$. Direct inversion gives
$P(X^{-1})=-X^{-s}P(X)$ and hence $Q(X^{-1})=X^{3-s}Q(X)$.
Choose an integer k with $2k\equiv s-3\pmod K$. Then
$\xi=\zeta^{-k}Q(\zeta)$ is real and is an algebraic integer.
The full cyclotomic norm of Q is the square of the real-subfield norm
of $\xi$, because $\operatorname{Norm}(\zeta)=1$.
Also $\operatorname{Norm}(\zeta-1)=K$. Therefore
\[
 \operatorname{Res}(\Phi_K,P)
       =K^3\operatorname{Norm}_{\mathbb Q(\zeta+\zeta^{-1})/\mathbb Q}(\xi)^2.
\]
The inner norm is an integer. In the non-forced case the additive rigidity
input implies $P(\zeta)\ne0$, so this integer is nonzero. This argument
uses reciprocity as a Laurent identity and does not assume that a leading
coefficient survives a finite-field specialization.
\end{proof}


The Fibonacci slice
\[
\Res(\Phi_K(X),X^2+3X+1)=F_K^2
\]
shows that individual collision resultants can be exponentially large.  At
$K=131$, $p=F_{131}$, a Lucas primality certificate and independent Python/C++
arbitrary-precision computations give exact nonzero Cayley excess $32$
(finite certificate:
\path{computations/verification/cayley_resultant_fibonacci.py} and
\path{computations/cpp/cayley_fibonacci_verifier.cpp}); this is
a finite certificate, not an asymptotic theorem.  The primitive-divisor input
for the associated rational model is from \cite{Yabuta2001}.

Let $S_K$ be the primitive squarefree nonzero-difference polynomial obtained
from the resultant of the characteristic-zero Cayley polynomial with its
translate.  For a split prime $p\equiv1\pmod K$, with $H$ and
$A=c(H\setminus\{1\})$ as in the preceding sections, define the
\emph{Cayley excess}
\begin{equation}\label{eq:Delta-def}
\Delta_K(p)
 =\sum_{d\ne0}r_{A-A}(d)^2-(K-1)(2K-5)
 =\sum_tC_H(t)^2-(K-1)(2K-5)\in\mathbf Z_{\ge0}.
\end{equation}
The two expressions agree by \eqref{eq:energyCH} and
$r_{A-A}(0)=|A|=K-1$.  Nonnegativity holds because the identity and central
inversion-swap collisions are distinct ordered collision quadruples already
contributing exactly $(K-1)(2K-5)$ to the first sum: $\Delta_K(p)$ counts
the remaining, non-forced, ordered collision quadruples.  Then for split
primes
\begin{equation}\label{eq:disc-class}
\Delta_K(p)>0\iff p\mid\Disc(S_K),
\qquad
\Delta_K(p)\le4v_p(\Disc(S_K)).
\end{equation}
The exact global form of $\Disc(S_K)$ comes from the ratio-orbit
factorization.  Write $n=K-1$ and $m=(K-1)/2$.  Then
$S_K=\prod_{[r]}F_r$, the product over the $m$ ratio classes
$[r]$, $r\sim r^{-1}$, where $F_r$ is the degree-$n$ Galois-orbit factor of
the differences $c(\zeta^a)-c(\zeta^{ra})$.  For one orbit, the index
formula in the maximal real subfield, whose discriminant is $K^{m-1}$,
gives $|\Disc(F_r)|=2^{n^2}K^{K-2}I_r^4$ for an integer $I_r>0$.  For two
distinct classes, $|\Res(F_r,F_s)|=2^{n^2}K^{K+1}M_{r,s}^2$; the square
comes from writing the even half-degree factors as $H_r(Y)=1+KP_r(Y)$ and
using $\deg(P_s-P_r)=m-1$.  That degree drop is the one place where
Dedekind sums enter: the coefficient of $Y^{m-1}$ in $P_r$ is
$(K-1)(K-2)/12-Ks(r,K)$, and by the prime-modulus equality criterion
\cite[Corollary~1.2]{JRW2011} (pp.~2197--2202), $s(r,K)=s(s,K)$ if and
only if $r\equiv s$ or $rs\equiv1\pmod K$, so distinct ratio classes have
distinct coefficients (for $K\ge5$; at $K=3$ there is a single ratio class
and the formula is a direct check).  Since
$\Disc(\prod_rF_r)=\prod_r\Disc(F_r)\prod_{r<s}\Res(F_r,F_s)^2$, collecting
the powers of $2$ and $K$ -- the exponents are
$mn^2+m(m-1)n^2=(K-1)^4/4$ and $m(K-2)+m(m-1)(K+1)=(K-1)(K^2-7)/4$ -- gives
\begin{equation}\label{eq:global-fourth}
|\Disc(S_K)|
 =2^{(K-1)^4/4}K^{(K-1)(K^2-7)/4}J_K^4
\end{equation}
for the integer
$J_K=\prod_{[r]}I_r\prod_{[r]<[s]}|M_{r,s}|>0$.  Exact resultant
computations through $K=47$ (the normalized prime-to-$2K$ mass
$4\log J_K^{(2K)}/K^4$ increases from $0.0837$ at $K=23$ to $0.1182$ at
$K=47$) show that the residual logarithmic mass genuinely remains
$K^4$-scale on the computed range, so \eqref{eq:global-fourth} is
structural, not a hidden proof of \eqref{eq:sharp-target}.

\section{Fixed-product rows and exact split support}

For one valid subgroup pair put $u=xy$, $s=x+y$.  Then
\begin{equation}\label{eq:rowparam}
c(x)+c(y)=\frac{2(u-1)}{u-s+1}.
\end{equation}
At fixed $u\ne1$, regard the parameter as the finite value
$\sigma=c(x)+c(y)$ in \eqref{eq:rowparam}; the case $u-s+1=0$ is the
excluded infinite parameter.  For finite $\sigma$ -- necessarily
$\sigma\ne0$, since $\sigma=0$ would force $u=1$ -- \eqref{eq:rowparam}
gives $s=u+1-2(u-1)/\sigma$, and $\{x,y\}$ is then the root pair of
$T^2-sT+u$; hence there is at most one unordered pair.
Let $I_K(p)$ be the maximum number of common parameters of two distinct product
rows.  The exact cross-row decomposition gives
\begin{equation}\label{eq:rowenergy}
\Delta_K(p)\le4I_K(p)(K-1)(K-2).
\end{equation}

Normalize a row pair to $(1,r)$ and set
\[
\tau=\zeta+\zeta^{-1},\qquad
\alpha_r=U_{r-1}(\tau/2),\qquad
\beta_r=-2U_{r-2}(\tau/2).
\]
The row equation is the affine trace relation
\begin{equation}\label{eq:affinetrace}
Y=\alpha_rX+\beta_r.
\end{equation}
Let $C_K(T)$ be the minimal polynomial of $\tau$ and set
\[
H_K(T)=(T-2)C_K(T).
\]
The two full row polynomials have one unique characteristic-zero common trace,
the invalid row trace.  Removing it produces fixed-degree source/target
cofactors over $\Z[\tau]$.  For every split embedding prime $\mathfrak q$,
\begin{equation}\label{eq:gcdcount}
\deg\gcd(P_{K,r},Q_{K,r}\bmod\mathfrak q)
 =I_{1,r}(\mathfrak q).
\end{equation}

Let the coefficient ideal of the first three fixed-degree subresultants
contract to $\Lambda_{K,r}^{(3)}\Z$.

\begin{theorem}[Exact real-trace order-three support]\label{thm:lambda}
For every split rational prime $p\equiv1\pmod K$,
\[
p\mid\Lambda_{K,r}^{(3)}
\]
if and only if some row pair in ratio class $[r]$ has at least three common
Cayley parameters.
\end{theorem}

\begin{proof}
Let $S_0,S_1,S_2$ be the first three fixed-degree subresultants of the two
cofactors, and let $J^{\mathrm{tr}}_{K,r}\subset\Z[\tau]$ be the ideal
generated by all their coefficients, so
$\Lambda_{K,r}^{(3)}\Z=J^{\mathrm{tr}}_{K,r}\cap\Z$.  This ideal is nonzero:
the cofactors have no common characteristic-zero trace (the unique common
trace was divided out), so $S_0\ne0$.

We use the contraction criterion: for a nonzero ideal $J$ of $\Z[\tau]$
with $J\cap\Z=\Lambda\Z$, a rational prime $p$ divides $\Lambda$ if and
only if $J\subseteq\mathfrak q$ for some prime ideal $\mathfrak q$ of
$\Z[\tau]$ above $p$.  Indeed, $J\subseteq\mathfrak q$ gives
$\Lambda\in\mathfrak q\cap\Z=p\Z$.  Conversely, if $p\mid\Lambda$ and
$J+p\Z[\tau]$ were the whole ring, localizing at $p$ would put
$1\in J_{(p)}$; clearing denominators in the domain $\Z[\tau]$ then gives
an element of $J\cap\Z$ prime to $p$, contradicting $p\mid\Lambda$.
Hence $J+p\Z[\tau]$ is proper, lies in a maximal ideal $\mathfrak q$, and
$\mathfrak q\cap\Z=p\Z$.

The valid cofactors remain fixed-degree and squarefree after split
reduction -- their leading coefficients are $1$ and the unit
$\alpha_r^{m+1}$, and $p\ne K$ -- so \eqref{eq:gcdcount} is exact and
there is no specialization degree drop.  By the criterion,
$p\mid\Lambda_{K,r}^{(3)}$ if and only if some $\mathfrak q\mid p$
contains every coefficient of $S_0,S_1,S_2$, i.e. if and only if
$S_0\equiv S_1\equiv S_2\equiv0\pmod{\mathfrak q}$.  Brown--Traub's
fixed-degree subresultant criterion \cite{BT1971} identifies this with
$\deg\gcd(P_{K,r},Q_{K,r}\bmod\mathfrak q)\ge3$, and \eqref{eq:gcdcount}
identifies that degree with $I_{1,r}(\mathfrak q)$.  As $\mathfrak q$
ranges over the primes above the split prime $p$, the associated real
embeddings cover all row pairs in ratio class $[r]$, which gives the
rational-prime statement.
\end{proof}

After removing factors $2,K$ and multiplying over ratio classes, write the
result as $\Lambda'_K{}^{(3)}$.  Then
\begin{equation}\label{eq:outside-support}
p\nmid\Lambda'_K{}^{(3)}
\Longrightarrow
I_K(p)\le2,
\quad
\Delta_K(p)\le8(K-1)(K-2),
\quad
E_+(A)\le11(K-1)(K-2).
\end{equation}
The older dehomogenized $D'/G'$ invariants remain valid only as one-way
obstructions because specialization can reduce leading degree at projective
infinity.

\section{Characteristic-zero affine-intercept rigidity}

Write
\[
t_a=\zeta^a+\zeta^{-a},\qquad
s_a=\zeta^a-\zeta^{-a},\qquad
\gamma_r(a,b)=t_b-\alpha_rt_a.
\]
Let $[a]$ denote the trace class $a\sim-a$, and put
$E_j=|1-\zeta^j|$, $1\le j\le m=(K-1)/2$.

\begin{lemma}[Prime-conductor multiplicative independence]\label{lem:bass}
The positive numbers $E_1,\ldots,E_m$ are multiplicatively independent.
\end{lemma}

\begin{proof}
Suppose $\prod_{j=1}^mE_j^{n_j}=1$ with $n_j\in\Z$.  Since
$E_j^2=(1-\zeta^j)(1-\zeta^{-j})$ and
$(1-\zeta^{-j})=-\zeta^{-j}(1-\zeta^j)$, squaring the relation gives
$\prod_j(1-\zeta^j)^{2n_j}=\omega$ for a root of unity $\omega$ (a sign
times a power of $\zeta$).  Taking norms from $\mathbf Q(\zeta)$: every
$1-\zeta^j$ is a Galois conjugate of $1-\zeta$, so
$N(1-\zeta^j)=N(1-\zeta)=\Phi_K(1)=K$, while $N(\omega)=1$.  Hence
$K^{2\sum_jn_j}=1$, forcing
\[
\sum_{j=1}^mn_j=0.
\]
The relation therefore rewrites as one among the ratios
$(1-\zeta^j)/(1-\zeta)$, $2\le j\le m$ (the $j=1$ ratio is $1$, and the
vanishing exponent sum absorbs the powers of $1-\zeta$).  By Bass's
theorem for prime conductor \cite{Bass1966} (Theorem~3), these ratios are
multiplicatively independent modulo torsion: their classes form a free
basis of the cyclotomic units modulo torsion.  Hence $2n_j=0$ for
$j\ge2$, and the vanishing sum gives $n_1=0$.
\end{proof}

Use the elementary identities
\begin{equation}\label{eq:absfacts}
|s_n|=E_{[2n]},\qquad
|t_x-t_y|=E_{[x-y]}E_{[x+y]}.
\end{equation}

\begin{theorem}[Trace-intercept fibre rigidity]\label{thm:intercept}
For every odd prime $K\ge7$ and every normalized ratio $r\ne1$, every fibre of
$([a],[b])\mapsto\gamma_r(a,b)$ has cardinality at most two.

If $r=-1$, a double fibre is only a swap.  If $r\ne\pm1$, every double fibre is,
up to signs and exchange,
\begin{equation}\label{eq:doublefibre}
([a],[a+1-r]),\qquad([a+2],[a+1+r]).
\end{equation}
\end{theorem}

\begin{proof}
Suppose two distinct pairs have the same intercept.  For $r\ne-1$,
\[
t_b-t_d=\alpha_r(t_a-t_c).
\]
Taking absolute values and applying \eqref{eq:absfacts} yields
\[
E_{[b-d]}E_{[b+d]}E_{[2]}
 =E_{[a-c]}E_{[a+c]}E_{[2r]}.
\]
Lemma~\ref{lem:bass} makes the multisets of folded indices equal.  Since
$[2r]\ne[2]$, one source difference must be $[2]$; after changing signs and
ordering take $c=a+2$.  The remaining two indices force the target classes in
\eqref{eq:doublefibre}.  The equality itself follows from
\[
t_{x+r}-t_{x-r}=s_rs_x,
\qquad
t_{a+2}-t_a=s_1s_{a+1}.
\]
A third point would form a triangle in the step-two quotient of the odd
$K$-cycle by $a\sim-a$; that quotient is a path.  For $r=-1$ the intercept is
$t_a+t_b$.  If $t_a+t_b=t_c+t_d$, subtracting gives a rational relation
among the powers $1,\zeta,\ldots,\zeta^{K-1}$ whose coefficient sum is
$4-4=0$.  The space of such relations is spanned by $(1,\ldots,1)$, the
coefficient vector of $\Phi_K$; the zero coefficient sum forces the scalar
to vanish, so every coefficient vanishes separately, the multisets
$\{\pm a,\pm b\}$ and $\{\pm c,\pm d\}$ in $\mathbf F_K$ agree, and
$\{[a],[b]\}=\{[c],[d]\}$: only identity or swap.  This covers in
particular $K=7$, where the relation has length $8>\deg\Phi_K=6$.
\end{proof}

The distinguished characteristic-zero row solution $([1],[r])$ is therefore
isolated.  This mechanism is related to cyclotomic-unit cross-ratio methods in
projective equivalence of roots of unity \cite{Fu2022}, but the affine theorem
above is proved directly in the prime-conductor setting.
\section{Cyclotomic-unit matching and a scalar shift resultant}

Put $u_n=s_n/s_1$.  For valid trace classes define
\begin{equation}\label{eq:WV}
W_r(b)=u_{(b+r)/2}u_{(b-r)/2},\qquad
V_r(a)=u_ru_{(a+1)/2}u_{(a-1)/2}.
\end{equation}
A direct difference factorization gives
\begin{equation}\label{eq:unitmatch}
\gamma_r(a,b)-\beta_r=s_1^2\bigl(W_r(b)-V_r(a)\bigr).
\end{equation}
Since $s_1$ is a unit away from $K$, every split row intersection is exactly a
congruence $W_r(b)\equiv V_r(a)\pmod{\mathfrak q}$ between explicit products
of cyclotomic units.

Let
\[
P_1(T)=\frac{H_K(T)}{T-t_1},\qquad
R_r(U)=\frac{H_K(U)}{U-t_r}
\]
and define the scalar shift resultant
\begin{equation}\label{eq:shiftres}
\mathcal R_{K,r}(Z)
 =\Res_T\!\left(P_1(T),R_r(\alpha_rT+\beta_r+Z)\right)
 =c_0+c_1Z+c_2Z^2+\cdots.
\end{equation}

\begin{theorem}[Scalar support identity]\label{thm:scalar-support}
For every split embedding prime $\mathfrak q$,
\begin{equation}\label{eq:ordshift}
\ord_{Z=0}(\mathcal R_{K,r}(Z)\bmod\mathfrak q)
 =I_{1,r}(\mathfrak q).
\end{equation}
Consequently
\[
I_{1,r}(\mathfrak q)\ge3
\iff
c_0\equiv c_1\equiv c_2\equiv0\pmod{\mathfrak q}.
\]
\end{theorem}

\begin{proof}
Over the split residue field the two valid trace polynomials factor linearly,
so the product formula gives
\[
\mathcal R_{K,r}(Z)
 =\prod_{[a]\ne[1]}\prod_{[b]\ne[r]}
   (Z+\alpha_rt_a+\beta_r-t_b).
\]
Each factor is linear in $Z$ and vanishes at zero exactly for one row
intersection.
\end{proof}

The ideal $(c_0,c_1,c_2)$ contracts to an integer
$\Omega_{K,r}^{(3)}$ with the same split rational-prime order-three support as
$\Lambda_{K,r}^{(3)}$.  Equality of the full integers is not asserted.
Moreover, the nonzero characteristic-zero constant coefficient gives the
per-ratio deterministic budget, where
$I_r(p):=\max_{\mathfrak q\mid p}I_{1,r}(\mathfrak q)$ over the embedding
primes $\mathfrak q$ of the split prime $p$:
\begin{equation}\label{eq:perratio}
\sum_{\substack{p\equiv1\,(K)\\I_r(p)\ge3}}\log p=O(K^3).
\end{equation}
Indeed, $c_0\ne0$ in characteristic zero: by
Theorem~\ref{thm:intercept} the distinguished intercept is isolated, so no
factor in the product formula for $\mathcal R_{K,r}(0)$ in
Theorem~\ref{thm:scalar-support} vanishes.  Every prime counted in
\eqref{eq:perratio} divides the nonzero integer
$|N_{\mathbf Q(\tau)/\mathbf Q}(c_0)|$: for such a prime, some
$\mathfrak q\mid p$ has $c_0\equiv0\pmod{\mathfrak q}$ by
Theorem~\ref{thm:scalar-support}, and $\mathfrak q$ has residue degree
one.  It remains to bound this norm.  Under the $j$-th real embedding,
$|\sigma_j(u_n)|=|\sin(2\pi jn/K)|/|\sin(2\pi j/K)|\le K/(4d_j)$ with
$d_j=\min(j,K/2-j)$, using $|\sin\pi x|\ge2\,\mathrm{dist}(x,\mathbf Z)$.
Since $\alpha_r=u_r$ and $\beta_r=-2u_{r-1}$, each of the $m^2$ factors of
$c_0$ has $j$-th embedding at most $2+K/d_j$ in absolute value, so
\[
\log|N(c_0)|
 \le m^2\sum_{j=1}^m\log(2+K/d_j)=O(K^3):
\]
the last sum is $O(K)$ by the elementary Stirling estimate applied to the
integers and half-integers $d_j$.  Summing $\log p$ over the distinct
counted primes, each dividing $|N(c_0)|$, gives \eqref{eq:perratio}.
Summing over $O(K)$ ratio classes recovers an $O(K^4)$ scale, so
\eqref{eq:perratio} is not a polynomial largest-prime theorem.

\section{Inverse reflection and the pair-sum square}

The ratio $r=-1$ has additional symmetry.  Here
\[
\alpha=-1,\qquad \beta=2\tau,
\qquad Q(T)=-P(2\tau-T).
\]
Set
\begin{equation}\label{eq:FEO}
F(X)=P(\tau+X)=E(X^2)+XO(X^2).
\end{equation}
The reflection reduction gives
\begin{equation}\label{eq:inverse-gcd}
I_{-1}(\mathfrak q)=2\deg\gcd(E,O\bmod\mathfrak q).
\end{equation}
Thus inverse-row intersections are automatically even.  A fixed-degree
half-size subresultant construction defines the exact split support integer
$\Xi_K$.

The scalar shift resultant admits a stronger universal factorization.

\begin{theorem}[Inverse pair-sum factorization]\label{thm:pairsum}
Let $x_1,\ldots,x_m$ be the roots of the monic polynomial $F$.  Define
\[
S(Z)=\prod_{i<j}(Z-x_i-x_j).
\]
Then $S\in\Z[\tau][Z]$ and
\begin{equation}\label{eq:pairsquare}
\Res_X(F(X),F(Z-X))=2^mF(Z/2)S(Z)^2.
\end{equation}
For every split embedding prime,
\begin{equation}\label{eq:inverse-order}
I_{-1}(\mathfrak q)=2\ord_{Z=0}(S(Z)\bmod\mathfrak q).
\end{equation}
Hence
\begin{equation}\label{eq:two-scalar}
I_{-1}(\mathfrak q)\ge3
\iff
S(0)\equiv S'(0)\equiv0\pmod{\mathfrak q}.
\end{equation}
\end{theorem}

\begin{proof}
The root-product formula gives
\[
\Res_X(F(X),F(Z-X))=\prod_{i,j}(Z-x_i-x_j).
\]
The diagonal ordered pairs give $2^mF(Z/2)$, and every unordered pair $i<j$
occurs twice.  The coefficients of $S$ are symmetric integral polynomials in
the roots.  At a split prime, $p\ne K$ is unramified and
$H_K\bmod\mathfrak q$ is separable -- its roots, the trace values $t_a$,
are pairwise distinct modulo $\mathfrak q$ -- so $\tau$ is a simple root
modulo $\mathfrak q$, i.e. $F(0)=P(\tau)=H'_K(\tau)\not\equiv0
\pmod{\mathfrak q}$ (and $\tau\not\equiv2\pmod{\mathfrak q}$).  Since $p$
is odd, the diagonal factor $2^mF(0)$ is a unit at $Z=0$.  Combining with the
general shift-resultant identity gives \eqref{eq:inverse-order};
\eqref{eq:two-scalar} follows because the left side is even.
\end{proof}

Let $S(Z)=s_0+s_1Z+\cdots$ and let the ideal $(s_0,s_1)$ contract to
$\Theta_K\Z$.  Then $\Theta_K$, $\Xi_K$, and the inverse factor of
$\Lambda_{K,-1}^{(3)}$ have the same split rational-prime support; their
valuations need not agree.

\section{The pair-sum polynomial as a translated half-degree resultant}

The square root $S$ itself can be obtained without constructing the
degree-$m^2$ shift resultant.  Translate
\begin{equation}\label{eq:translatedF}
F_Z(X)=F(X+Z/2)=E_Z(X^2)+XO_Z(X^2).
\end{equation}

\begin{theorem}[Translated parity-resultant identity]\label{thm:translated-res}
There exists a convention-dependent sign $\varepsilon_m\in\{\pm1\}$ such that
\begin{equation}\label{eq:translated-res}
S(Z)=\varepsilon_m\Res_Y(E_Z(Y),O_Z(Y)).
\end{equation}
Consequently
\begin{align}
S(0)&=\varepsilon_m\Res_Y(E,O),\label{eq:s0res}\\
S'(0)&=\varepsilon_m
\left.\frac{d}{dZ}\Res_Y(E_Z,O_Z)\right|_{Z=0}.\label{eq:s1res}
\end{align}
The first-order deformation is
\begin{equation}\label{eq:deformation}
\left.\partial_ZE_Z\right|_0=\frac{O+2YO'}{2},
\qquad
\left.\partial_ZO_Z\right|_0=E'.
\end{equation}
\end{theorem}

\begin{proof}
The roots of $F_Z$ are $y_i=x_i-Z/2$.  Apply the constant pair-sum identity
from Theorem~\ref{thm:pairsum} to the monic polynomial $F_Z$.  Its unordered
pair-sum product at zero is
\[
\prod_{i<j}(-y_i-y_j)
 =\prod_{i<j}(Z-x_i-x_j)=S(Z).
\]
It remains to identify this product with the parity resultant.  For a
monic polynomial $G$ with parity decomposition
$G(X)=\mathcal E(X^2)+X\mathcal O(X^2)$ and roots $y_i$, the constant
pair-sum coefficient satisfies $S_G(0)^2=\Res_Y(\mathcal E,\mathcal O)^2$.
Indeed, \eqref{eq:pairsquare} at $Z=0$ gives
$\Res_X(G(X),G(-X))=2^mG(0)S_G(0)^2$, while the root-product formula and
$G(-y_i)-G(y_i)=-2y_i\mathcal O(y_i^2)$ give
$\Res_X(G(X),G(-X))=\prod_iG(-y_i)
=2^m\mathcal E(0)\Res_Y(\mathcal E,\mathcal O)^2$:
here $\prod_iy_i=(-1)^m\mathcal E(0)$, and
$\prod_i\mathcal O(y_i^2)=\Res_X(G,\mathcal O(X^2))
=\Res_Y(\mathcal E,\mathcal O)^2$, evaluating at the roots $\pm\sqrt y$ of
$\mathcal O(X^2)$ and using that exactly one of $\mathcal E,\mathcal O$
is monic, in either parity case.  Since $G(0)=\mathcal E(0)\ne0$ in
characteristic zero, $S_G(0)^2=\Res_Y(\mathcal E,\mathcal O)^2$.  Apply
this to $G=F_Z$ over $\mathbf Q(\tau)(Z)$: its pair-sum product at zero
is $S(Z)$ by the display above, and its parity components are $E_Z,O_Z$
by \eqref{eq:translatedF}.  Hence
$S(Z)^2=\Res_Y(E_Z,O_Z)^2$ as polynomials in $Z$; both are nonzero in
the domain $\Z[\tau][Z]$, so they differ by a constant sign
$\varepsilon_m\in\{\pm1\}$, fixed by the monicity of $S$ once the
resultant convention is fixed.  Differentiating
$F(X+Z/2)$ at $Z=0$ gives \eqref{eq:deformation}.
\end{proof}

Thus inverse order-three support is the simultaneous divisibility problem for a
half-degree resultant and its translation derivative.  In exact computation,
both scalars can be obtained in one Sylvester determinant over dual numbers
$\varepsilon^2=0$, reducing the determinant dimension to $O(K)$ rather than
materializing a degree-$m^2$ polynomial.

Exact Smith contractions give, after removing factors $2K$,
\[
\Theta'_7=1,\quad
\Theta'_{11}=1,\quad
\Theta'_{13}=5,\quad
\Theta'_{17}=1,\quad
\Theta'_{19}=37,\quad
\Theta'_{23}=47^2\cdot137.
\]
For comparison $\Xi'_{23}=47\cdot137$: split prime support agrees, valuations
differ.  The reference SymPy implementation reduces the $K=17$ computation
from roughly $51$ seconds for the full shift resultant to roughly $0.6$ seconds
for the translated half-degree route; at $K=23$ the latter remains on the
order of seconds, after which Smith normal form becomes the principal finite
algebra bottleneck.  These timings are reproducibility data, not asymptotic
complexity theorems.

The norm of the inverse constant shift coefficient also has the exact shape
\begin{equation}\label{eq:inverse-norm}
\left|\Norm(R_-(0))\right|
 =2^{m^2}K^m\left|\Norm(\Res(E,O))\right|^2,
\end{equation}
using $|\Norm(F(0))|=K^m$.  Hence every prime away from $2K$ occurs with even
valuation in this particular norm.  This explains part of the observed square
pattern but does not imply equality of the support integers.

\section{Global split-prime budget and minimal-extension Frobenius}

Simultaneous multiplication of collision exponent indices by
$a\in\mathbf F_K^\times$ acts freely on non-forced collision quadruples.  The
number of normalized scaling orbits is
\[
N_K=K^3-5K^2+6K+1.
\]
Combining the square-resultant theorem with $|\Res(\Phi_K,P)|\le12^{K-1}$
gives:

\begin{theorem}[Weighted split-prime budget]\label{thm:budget}
\begin{equation}\label{eq:budget}
\sum_{p\equiv1\,(K)}\Delta_K(p)\log p
\le N_K\bigl((K-1)\log12-3\log K\bigr).
\end{equation}
\end{theorem}

\begin{proof}
Index the elements of $A$ by $\mathbf F_K^\ast$.  An admissible
collision quadruple is $q=(i,j,r,s)$ with $i\ne j$, $r\ne s$;
simultaneous multiplication by $a\in\mathbf F_K^\ast$ acts freely,
every orbit has a unique representative with $i=1$, and the number
of non-forced orbits is
\[
N_K=\frac{(K-1)^2(K-2)^2-(K-1)(2K-5)}{K-1}
   =(K-1)(K-2)^2-(2K-5)=K^3-5K^2+6K+1.
\]
For a normalized representative let $P_q$ be the cleared Cayley
collision polynomial and $R_q=\Res(\Phi_K,P_q)$.  By the
square-resultant theorem $R_q=K^3M_q^2$ with $M_q\ne0$, and
$|R_q|\le12^{K-1}$.  For a split prime $p$ and a generator $h$ of
$H$, the equation $P_q(h^a)=0$ is exactly the collision attached to
$a\cdot q$; writing $s_q(p)$ for the number of such $a$, one has
$\Delta_K(p)=\sum_q s_q(p)$.  Since $\Phi_K$ is separable modulo
$p$ and splits completely, every root counted by $s_q(p)$
contributes a factor $p$ to $R_q$, hence
$s_q(p)\le v_p(R_q)=2v_p(M_q)$, the last equality because $p\ne K$.
Summing over split primes and using
$\sum_p v_p(M_q)\log p\le\log|M_q|$ gives
\[
\sum_p\Delta_K(p)\log p
\le\sum_q2\log|M_q|
=\sum_q\log(|R_q|/K^3)
\le N_K\bigl((K-1)\log12-3\log K\bigr).
\]
\end{proof}

For prime $K\asymp X^\theta$ and split prime $p\asymp X$,
Bombieri--Vinogradov and PNT \cite{Bombieri1965,IK2004} imply
\[
\operatorname{average}\frac{\Delta_K(p)}{K^2}
 \ll_\theta X^{3\theta-1}
\]
for fixed $\theta<1/3$.  Integrality of the excess strengthens this to a
density-one exact baseline for fixed $\theta<1/5$:
\[
\Delta_K(p)=0,
\qquad
E_+(A)=3(K-1)(K-2)
\]
for asymptotically almost every split pair.
Indeed, the average is over the dyadic family $Q/2<K\le Q$,
$X/2<p\le X$, $p\equiv1\pmod K$, whose cardinality has order
$X/(\log X)^2$ by Bombieri--Vinogradov and the prime number
theorem.  Since $\Delta_K(p)$ is a nonnegative integer and
$K^2\le Q^2=X^{2\theta}$ on this family, Markov's inequality gives
\[
\frac{\#\{(K,p):\Delta_K(p)\ge1\}}{|\mathcal D_\theta(X)|}
\le\operatorname{average}\Delta_K(p)
\le Q^2\cdot\operatorname{average}\frac{\Delta_K(p)}{K^2}
\ll_\theta X^{5\theta-1}\longrightarrow0
\qquad(\theta<1/5),
\]
and on the complementary pairs $\Delta_K(p)=0$, whence
$E_+(A)=(K-1)^2+(K-1)(2K-5)=3(K-1)(K-2)$.

Now let $p\ne2,K$ be arbitrary and put
\[
f=\ord_K(p),\qquad q=p^f.
\]
In $\mathbf F_q$ the order-$K$ subgroup exists uniquely.  Put
$A_{K,p}=c(H\setminus\{1\})\subset\mathbf F_q$ and
\[
\Delta_K^{\rm ext}(p)
 =\sum_{d\ne0}r_{A_{K,p}}(d)^2-(K-1)(2K-5),
\]
the number of ordered non-forced collision quadruples in the
minimal extension.  It is a nonnegative integer, and
$\Delta_K^{\rm ext}(p)=0$ gives $E_+(A_{K,p})=3(K-1)(K-2)$.

\begin{theorem}[Frobenius amplification]\label{thm:frob}
Every non-forced collision orbit under the $p$-power Frobenius has size exactly
$f$.  Hence
\[
f\mid\Delta_K^{\rm ext}(p).
\]
The weighted resultant budget and the discriminant classification extend to
all such rational primes in the minimal extension.
\end{theorem}

\begin{proof}
Frobenius multiplies every exponent by $p$ modulo $K$.  If a positive power
$d<f$ fixed a collision tuple, any nonzero coordinate would imply
$p^d\equiv1\pmod K$, contradicting the definition of $f$.  Hence every
non-forced Frobenius orbit has size exactly $f$, proving
$f\mid\Delta_K^{\rm ext}(p)$.

It remains to extend the budget and the classification.  Retain the
$N_K$ normalized scaling orbits and their resultants
$R_Q=\Res(\Phi_K,P_Q)=K^3M_Q^2$ from the split budget.  For $p\ne K$
choose a prime ideal $\mathfrak P$ of $\mathbf Z[\zeta_K]$ above
$p$; its residue degree is $f$, and reduction sends $\zeta_K$ to a
primitive $K$th root $h\in\mathbf F_q$.  With
$s_Q(p)=\#\{a\in\mathbf F_K^\ast:P_Q(h^a)=0\}$ one has
$\Delta_K^{\rm ext}(p)=\sum_Qs_Q(p)$, and the product formula
$R_Q=\prod_aP_Q(\zeta_K^a)$ puts $R_Q\in\mathfrak P^{s_Q(p)}$.
Since $p\ne K$ is unramified, $v_{\mathfrak P}$ of a rational
integer is the ordinary $p$-valuation, so
$s_Q(p)\le v_p(R_Q)=2v_p(M_Q)$.  Summing over $p$ exactly as in the
split case gives
$\sum_{p\ne2,K}\Delta_K^{\rm ext}(p)\log p
 \le N_K\bigl((K-1)\log12-3\log K\bigr)$.
For the classification, apply the maintained cluster argument to
the squarefree difference polynomial $S_K$ at $\mathfrak P$: its
roots are $\mathfrak P$-integral, distinct roots coalesce modulo
$\mathfrak P$ exactly when collisions occur, and
$v_{\mathfrak P}(\Disc S_K)=v_p(\Disc S_K)$; hence
$\Delta_K^{\rm ext}(p)>0$ if and only if $p\mid\Disc(S_K)$, with
$\Delta_K^{\rm ext}(p)\le4v_p(\Disc(S_K))$.
\end{proof}

PNT, Brun--Titchmarsh and the scarcity of small multiplicative orders then give
a density-one exact baseline for unrestricted prime pairs
$K\asymp X^\theta$, $p\asymp X$ for every fixed $\theta<1/3$.  This mechanism
is absent on the split slice because there $f=1$.  In detail, split
the pairs according as $f=\ord_K(p)\le K^\beta$ or $f>K^\beta$,
where $\max(0,4-1/\theta)<\beta<1$; such a $\beta$ exists exactly
when $\theta<1/3$.  The residue classes modulo $K$ of order at
most $K^\beta$ number at most
$\sum_{d\mid K-1,\,d\le K^\beta}\varphi(d)\le K^\beta\tau(K-1)
 =K^{\beta+o(1)}$,
and Brun--Titchmarsh bounds the primes in each such class in
$(X/2,X]$ by $O_\theta(X/(K\log X))$; summing over prime
$K\asymp Q$ makes the small-order pairs a fraction
$O(Q^{\beta-1+o(1)})$ of all pairs.  If $f>K^\beta$ and the pair is
bad, the orbit theorem gives
$\Delta_K^{\rm ext}(p)\ge f>K^\beta$, so the extended budget bounds
the bad primes for each $K$ by $O(K^{4-\beta}/\log X)$, a fraction
$O(Q^{4-\beta}/X)$ of all pairs.  Both fractions tend to zero.

\section{Full Cayley--Jacobi trace and corrected off/off geometry}

Let $S=H^\perp$ be the multiplicative-character annihilator and define the
normalized twisted-Jacobi matrix
\[
M_S(q)=\frac{e_p(-q/2)}{p-1}J_S(q).
\]
Character orthogonality gives the exact trace identity
\begin{equation}\label{eq:jacobi}
\widehat{1_A}(r)=K\operatorname{Tr}M_S(2r),
\end{equation}
and therefore
\begin{align}
\sum_q|\operatorname{Tr}M_S(q)|^2&=\frac{p(K-1)}{K^2},\label{eq:trace2}\\
\sum_q|\operatorname{Tr}M_S(q)|^4&=\frac{pE_+(A)}{K^4}.\label{eq:trace4}
\end{align}
The normalized fourth/second moment ratio tends to $3$ in the maintained split
double average below $\theta<1/3$.  This is a fourth-moment theorem, not a CLT.

The pure off/off family contains a coherent reducible fibre.  For generic
$\lambda=(h_1,h_2)$ it occurs at
\[
t_\lambda=\frac{2}{c(h_1)-c(h_2)}
\]
and has exactly $2p-8$ valid affine points.  After extracting the spike,
Igusa's arithmetic open-image theorem \cite{BLV2009} gives geometric
irreducibility: the constants quotient is procyclic, so the geometric image
contains the commutator of an open arithmetic image, hence an open subgroup
of $\mathrm{SL}_2(\mathbf Z_\ell)$ for $\ell\ne p$.
The normalized elliptic local systems are pure of weight zero on their smooth
lisse loci; their multiplicative middle-extension stalks instead have
normalized eigenvalues $\pm p^{-1/2}$.
Apply the Fouvry--Kowalski--Michel trace-correlation argument \cite{FKM2013}
on the common smooth open and restore the boundedly many omitted points using
their bounded normalized traces.  This gives square-root quasi-orthogonality
between non-isomorphic geometric parameter classes.

\begin{theorem}[Corrected off/off variance]\label{thm:offoff}
For every prime $p>3$ and every subgroup $H\le\mathbf F_p^\ast$ of
odd order $K$,
\begin{equation}\label{eq:offoff}
\Var_{t\ne0}O(t)
\ll pE_+(A)+K^2p+K^4p^{1/2}+K^4,
\end{equation}
with an absolute implied constant.
\end{theorem}

\begin{proof}
Write $O(t)=\sum_{\lambda\in H^2}O_\lambda(t)$ and use Cayley
coordinates $a=c(u)$, $b=c(v)$, $\alpha=c(h_1)$, $\beta=c(h_2)$,
$s=\alpha+\beta$, $d=\alpha-\beta$.  On the generic set
$\Lambda_g=\{h_1,h_2\ne1,\ h_1\ne h_2,\ h_1h_2\ne1\}$, of size
$(K-1)(K-3)$, clearing denominators in
$2/(a-b)+2/(T_\alpha(a)-T_\beta(b))=t$ gives a branch quartic whose
discriminant factors with residual quartic $R_{(s,d)}(t)$; one has
$\Disc_tR_{(s,d)}
 =2^{16}d^2s^2(s-d-2)^2(s-d+2)^2(s+d-2)^2(s+d+2)^2\ne0$
on $\Lambda_g$ and $R_{(s,d)}(2/d)=16/d^2$.  Hence the generic
pencil has five distinct finite bad fibres: $t_\lambda=2/d$ of
multiplicative type $I_4$ and four of type $I_1$.  At $t_\lambda$
the cleared cubic factors into the two distinct lines
$(\beta-1)(a-1)-(\alpha-1)(b-1)$ and
$(\beta+1)(a+1)-(\alpha+1)(b+1)$; removing the seven distinct
invalid points from their $2p-1$ affine points gives the exact
spike count $O_\lambda(t_\lambda)=2p-8$.  With $a_\lambda(t)$ the
middle-extension Frobenius trace of the Jacobian pencil,
\[
O_\lambda(t)
 =p+p\,\mathbf 1_{t=t_\lambda}-a_\lambda(t)+O(1)
\]
uniformly in $\lambda$ and $t$.  The pencil is non-isotrivial, so
by Igusa's open-image theorem \cite{BLV2009} the geometric image
contains the commutator of an open arithmetic image, hence an open
subgroup of $\mathrm{SL}_2(\mathbf Z_\ell)$: the rank-two local
system is geometrically irreducible, of bounded conductor (at most
six tame singular places).  Equal singular support forces
$(s',d')=(s,d)$ or $(-s,d)$, so each geometric isomorphism class
contains at most the two parameter classes $\lambda$ and
$\lambda^\ast=(h_2^{-1},h_1^{-1})$.  For non-isomorphic classes,
the trace-correlation argument of \cite{FKM2013} --- run on the
common smooth open via Schur orthogonality, the
Grothendieck--Lefschetz formula, Deligne's Riemann hypothesis and
the Grothendieck--Ogg--Shafarevich bound, with the boundedly many
omitted points restored at cost $O(1)$ in the normalized sum ---
gives $\sum_t a_\lambda(t)\overline{a_\mu(t)}=O(p^{3/2})$; the
$O(K^2)$ same/involution ordered pairs contribute $O(p^2)$ each by
the Weil bound, so
$\sum_t\bigl|\sum_{\lambda\in\Lambda_g}a_\lambda(t)\bigr|^2
 =O(K^2p^2+K^4p^{3/2})$.
The spike multiplicity $D(t)=\#\{\lambda\in\Lambda_g:
t_\lambda=t\}$ satisfies $\sum_{t\ne0}D(t)^2\le E_+(A)$, since
$t\mapsto2/t$ identifies it with Cayley difference counts.  The
complement $H^2\setminus\Lambda_g$ has $4K-3$ classes, each of
variance $O(p)$ (a bilinear count for $h_1=h_2$, Hasse bounds for
$h_1h_2=1$ and for exactly one $h_i=1$, and the zero class at
$(1,1)$), hence total variance $O(K^2p)$ by Cauchy--Schwarz.
Combining the generic and special variances gives
\eqref{eq:offoff}.
\end{proof}

Thus the off/off layer has collapsed onto the same pointwise Cayley-energy
bottleneck as the mixed layer.

\section{Exact trace-to-local-statistics bridge}

Let $U$ be an $N\times N$ unitary matrix with eigenvalues
$e^{2\pi i\theta_1},\ldots,e^{2\pi i\theta_N}$.  For a zero-sum integer tuple
$m_1+\cdots+m_k=0$, define
\[
D_U(m_1,\ldots,m_k)
 =\sum_{j_1,\ldots,j_k\ \mathrm{distinct}}
 e^{2\pi i\sum_{a=1}^km_a\theta_{j_a}}.
\]
For a set partition $\pi$ of $\{1,\ldots,k\}$, put
\[
\mu(\pi)=\prod_{B\in\pi}(-1)^{|B|-1}(|B|-1)!.
\]

\begin{theorem}[Factorial trace identity]\label{thm:factorial}
\begin{equation}\label{eq:factorial}
D_U(m_1,\ldots,m_k)
 =\sum_{\pi}\mu(\pi)
   \prod_{B\in\pi}
   \operatorname{Tr}U^{\sum_{a\in B}m_a}.
\end{equation}
\end{theorem}

This is M\"obius inversion on the partition lattice.  Let
$\phi\in\mathcal S(\mathbf R^{k-1})$ and periodize it at the mean-spacing scale.
Poisson summation gives the exact microscopic statistic
\begin{equation}\label{eq:trace-local}
C_{k,N}(U;\phi)
 =N^{-k}\sum_{m\in\Z^{k-1}}
  \widehat\phi(m/N)D_U(m_1,\ldots,m_k),
\end{equation}
where $m_k=-(m_1+\cdots+m_{k-1})$.  In particular,
\begin{equation}\label{eq:pair-local}
C_{2,N}(U;\phi)
 =N^{-2}\sum_{m\in\Z}\widehat\phi(m/N)
  \bigl(|\operatorname{Tr}U^m|^2-N\bigr).
\end{equation}
If $\widehat\phi$ has compact support, the relevant powers satisfy
$|m|=O(N)$: microscopic local statistics probe the Heisenberg scale.

A weighted comparison of the quantities $D_U(m)$ with the matching random
matrix ensemble at every compact Fourier support is sufficient for the
corresponding band-limited microscopic correlations. Under positive normalized
averaging, the next theorem removes the band-limit. The required joint-trace
estimates for the correct desymmetrized cat-map block remain open.

\subsection{Positive correlation measures remove the band-limit}
Write $d=k-1$, unfold the circle to length $N$, and let $x_j=N\theta_j$.
The lifted anchored factorial measure is
\[
 \mu_N=\frac1N\sum_{j_1,\ldots,j_k\ \mathrm{distinct}}
 \sum_{\ell\in\mathbb Z^d}
 \delta_{(x_{j_1}-x_{j_k}+N\ell_1,\ldots,x_{j_d}-x_{j_k}+N\ell_d)}.
\]
It is positive, even when eigenvalues coincide, and
$\mu_N(\phi)=C_{k,N}(U;\phi)$.

\begin{lemma}[Uniform translated-cube bound]
\label{lem:factorial-cube-bound}
For $N>2$ and every $a\in\mathbb R^d$,
\[
 \mu_N(a+[0,1)^d)\le 2k^k\bigl(1+\mu_N([-2,2]^d)\bigr).
\]
The same bound holds under positive normalized averaging.
\end{lemma}
\begin{proof}
Let $n(t)$ count the periodically lifted points in $[t,t+2)$, with
multiplicity, and put $a_k=0$. For each tuple in the cube, integrate starting
points $t$ in the unit interval immediately preceding its anchor. Each of its
other points lies in $[t+a_i,t+a_i+2)$ apart from irrelevant endpoints.
Since $2<N$, every window contains at most one lift of each labelled point.
Integration over a period and then H\"older's inequality give
\[
 \mu_N(a+[0,1)^d)
 \le\frac1N\int_0^N\prod_{i=1}^k n(t+a_i)\,dt
 \le\frac1N\int_0^N n(t)^k\,dt.
\]
For the falling factorial $(r)_k$ one has
$r^k\le k^k((r)_k+r)$ for nonnegative integers $r$.
Indeed, for $r<k$ use $r^k\le k^{k-1}r$; for $r\ge k$ each factor in
$(r)_k$ is at least $r/k$. Moreover,
\[
 \frac1N\int_0^N n(t)\,dt=2,
\]
\[
 \frac1N\int_0^N(n(t))_k\,dt
 =\int\bigl(2-\operatorname{span}(0,z_1,\ldots,z_d)\bigr)_+\,d\mu_N(z)
 \le2\mu_N([-2,2]^d).
\]
The identity counts the possible starts of a common length-two window for
each distinct-label tuple; their total length is the displayed positive part.
This proves the bound. Positivity and normalization preserve it on averaging.
\end{proof}

\begin{theorem}[Extension from all compact Fourier supports]
\label{thm:factorial-schwartz-extension}
Fix $k\ge2$. Let $\mu_N,\nu_N$ be positive normalized averages of lifted
$k$-point factorial measures of $N$-point spectra. Suppose that $\nu_N$ has a
limit on all Schwartz tests and
\[
 \mu_N(f)-\nu_N(f)\longrightarrow0
\]
for every Schwartz $f$ with compact Fourier support, with no common upper
bound on that support. Then the same comparison holds for every Schwartz $f$.
\end{theorem}
\begin{proof}
Choose a nonzero nonnegative smooth compactly supported Fourier bump $\eta$,
and put $h=\mathcal F^{-1}\eta$. Since $h(0)>0$, for sufficiently small
$\epsilon>0$ the nonnegative band-limited Schwartz function
\[
 B(z)=4|h(\epsilon z)|^2/h(0)^2
\]
majorizes one on $[-2,2]^d$. The hypotheses bound $\nu_N(B)$ and
$\mu_N(B)$. Lemma~\ref{lem:factorial-cube-bound} therefore bounds both
measures on every translated unit cube by one constant $M$.
For $p>d$ and every Schwartz $g$ this implies
\[
 |\mu_N(g)|+|\nu_N(g)|
 \le C_{d,p}M\sup_z(1+|z|)^p|g(z)|.
\]
For a smooth compactly supported $\chi$ equal to one near zero, set
$f_R=\mathcal F^{-1}(\chi(\cdot/R)\widehat f)$.
Then $f_R\to f$ in Schwartz topology. The last inequality makes the two
approximation errors uniform in $N$. First send $N\to\infty$ at fixed $R$,
then $R\to\infty$.
\end{proof}

Thus a separate absolute high-frequency trace-tail estimate is unnecessary
for this positive-average, all-bandwidth criterion. This does not establish
such an absolute estimate, or extend a comparison confined to one fixed cap.
For example, for even $N$, the unit clock and a period-two configuration with
basis $\{0,a\}$, $0<a<2$, $a\ne1$, have pair measures
\[
 \sum_{r\in\mathbb Z}\delta_r-\sum_{r\in\mathbb Z}\delta_{Nr},
 \qquad
 \sum_{r\in\mathbb Z}\delta_{2r}
 +\frac12\sum_{r\in\mathbb Z}(\delta_{2r+a}+\delta_{2r-a})
 -\sum_{r\in\mathbb Z}\delta_{Nr}.
\]
They agree on tests supported in Fourier space inside $(-1/2,1/2)$ but
have different full limits. The theorem also requires positivity of the
averages. It addresses correlation tests, not a spacing-distribution theorem
without the additional point-process hypotheses needed for that conclusion.
